\chapter{LITERATURE REVIEW}

\section{Evolution of Terrain Representation} Terrain representation has evolved
significantly over the last three decades. Early systems relied on manual
surveying and simple topographic maps. These methods lacked the detail needed
for automated navigation.

In 2000, the Shuttle Radar Topography Mission (SRTM) provided the first
near-global elevation dataset \cite{Franks2023}. Shortly after, the ASTER Global Digital Elevation
Model (GDEM) was released in 2009 \cite{Pan2020}. These legacy datasets provided crucial
elevation grids. However, they exhibited significant errors in high-altitude
areas \cite{Franks2023}. Steep cliffs and deep valleys caused radar distortions \cite{Pan2020, Franks2023}. These distortions
led to inaccurate peak heights \cite{Pan2020}.

To address these errors, researchers evaluated newer satellite datasets.
In 2020, the European Space Agency released the Copernicus GLO-30 model. Recent
evaluations confirm its superior vertical accuracy \cite{Franks2023}. It
out-performs older models in steep mountain terrain \cite{Franks2023}. At 30-meter resolution,
this model captures the unique shape of Himalayan peaks. The resulting database
is small enough to run on portable devices. This makes it an ideal foundation
for modern terrain-based positioning systems \cite{Pan2020, Steger2022}.

\section{Developments in Skyline Segmentation} Extracting a clear skyline from a
photograph is a well-studied problem. Early approaches in the 2000s used classic
edge-detection filters. These filters detected sudden changes in image
intensity.

However, traditional edge-detection failed in real-world scenarios. Haze,
clouds, and shadows often obscured the skyline. This led to broken or incorrect
boundaries \cite{Pan2020}.

To solve this issue, researchers transitioned to machine learning. In 2015,
convolutional neural networks revolutionized image segmentation. These deep
models could identify the sky even in poor lighting conditions.

However, early neural networks were too heavy for mobile phones. In response,
researchers developed compact architectures. MobileUNET became a popular choice
for mobile devices. It simplifies calculations to save battery power while
maintaining high pixel accuracy.

Other modern models adapt object detection frameworks. For example, YUNet adapts
the YOLOv11 framework for rapid sky detection \cite{Yang2025}. It achieves high
segmentation scores under difficult weather conditions \cite{Yang2025}.

A persistent challenge was the lack of labeled training images. In 2017, Brejcha
and Cadik introduced the GeoPose3K dataset \cite{Brejcha2017}. This library
contains over 3,000 precisely labeled mountain landscapes \cite{Brejcha2017}. In 2022, Tomesek et
al. explored cross-modal rendering to generate more training data
\cite{CrossLocate}. These datasets allow modern networks to learn real-world
features successfully.

\section{Progress in Silhouette Matching and Heading Resolution} Once the
skyline is extracted, it must be matched against the terrain database. This step
calculates the camera's location and direction \cite{Pan2020}.

Historically, matching relied on simple pixel-by-pixel comparisons. This method
was slow and easily confused by different camera lenses. It also required prior
knowledge of the camera's direction.

To resolve this, researchers implemented a rotational sliding-window search \cite{Pan2020}. The
system compares the query photo against a 360-degree virtual horizon at multiple
positions \cite{Pan2020}.

To make this process run in real time, researchers developed two-stage search
systems. In 2013, Rakthanmanon et al. demonstrated fast subsequence matching
under Dynamic Time Warping (DTW) \cite{Rakthanmanon2013}. DTW is a robust
shape-alignment technique. It adjusts for perspective distortions caused by
different camera lenses \cite{Rakthanmanon2013}.

Running DTW across a large database remains computationally expensive \cite{Rakthanmanon2013}. To solve
this bottleneck, Mueen et al. introduced a fast search algorithm in 2022
\cite{Mueen2022}. This algorithm performs rapid pruning by computing simple
distance profiles \cite{Mueen2022}. It quickly filters out bad candidates. This leaves only a few
profiles for DTW alignment. This combination enables fast, real-time location
matching on mobile hardware.

\section{Justification for GNSS-Based Validation} Developing an alternative
positioning system requires a reliable testing standard. Researchers
traditionally use high-accuracy GPS data to validate their systems.

Early positioning studies relied on small, custom photo sets. These sets lacked
diverse lighting conditions and precise coordinates.

This issue was resolved with the release of the Google Street View Trekker
dataset \cite{Everest_StreetView}. In 2014, researchers captured over 45,000
high-resolution panoramas along the Everest trail \cite{Everest_StreetView}. The Trekker hardware uses a
multi-lens camera array. It also integrates precise satellite receivers and
motion sensors. This hardware ensures highly accurate location tags for every
image \cite{Everest_StreetView}.

This georeferenced dataset provides an ideal reference map for our project. It
contains thousands of real-world mountain views with verified coordinates. We
use this dataset to measure our system's accuracy in the actual Khumbu
environment.
